\section{Problem Statement}

% Owner-authored in substance. Transcribed from drafts/problem-statement-v3.md, the
% finalized statement, into the paper's definition and assumption environments.
% No content added beyond that file. The setting prose descends unchanged from
% version 2 (DEC-007); the three layers and the coverage question come from DEC-015.
% See AI/DECISIONS.md for the basis of each settled value, and AI/STATUS.md for the
% items still open. The coverage target appears as a parameter because DEC-015
% records its value as unsettled. The Fithian et al. result in the final paragraph
% carries no \citet call and no bibliography entry, because paper/references.bib
% forbids an entry for a source nobody here has opened. A background agent read it;
% the owner has not. It needs an SRC entry and a bib entry before submission.

A social app lets users post a check-in at a public venue. A check-in records one drink, so one
visit can produce several check-ins. Each check-in carries a user, a timestamp, and a venue. The
app accepts a check-in only at the venue's own location. It serves a community of young adults
who want to socialize with friends. We treat the community as a finite set, fixed only at the
moment the app judges it, and its size can differ from one night to the next.

We distinguish three objects. The community network carries the friendships that decide who goes
out together. The adopter set carries the members who use the app. The app's own graph carries
the follows it has recorded. Only the third is observable.

\begin{definition}[The three layers]
Let $C = (V, E)$ be the \emph{community network}, where $V$ is the community and $E$ holds the
mutual friendships among its members. Let $A \subseteq V$ be the \emph{adopter set}, the members
who use the app, and call $|A| / |V|$ the \emph{adoption rate}. Let $G = (A, F)$ be the
\emph{app graph}, where $F \subseteq \{\, e \in E : e \subseteq A \,\}$ holds the mutual follows
the app has recorded. The app observes $G$. It observes neither $C$ nor $V \setminus A$.
\end{definition}

\begin{assumption}[No non-adopter signal]
Members of $V \setminus A$ post nothing.
\end{assumption}

\begin{assumption}[Adopter- and venue-dependent posting probability]
A present adopter posts at least once with positive probability. That probability can depend on
the adopter and on the venue.
\end{assumption}

\begin{assumption}[Friends at the same venue and time]
Non-friends post independently of one another. Friends at the same venue and time do not.
\end{assumption}

Assumption 3 quantifies dependence over $E$, not over $F$. The app therefore faces correlation
whose source it need not have recorded. Assumption 1 bounds the consequence. Correlated posting
arises only between two co-present adopters joined in $E$, since non-adopters post nothing and
non-friends post independently. The following assumption makes that bound exact.

\begin{assumption}[Information-complete app graph]
$F = \{\, e \in E : e \subseteq A \,\}$.
\end{assumption}

Assumption 4 does not claim that the app records friendships perfectly. It states the case in
which the app holds every tie capable of producing a correlated post. We take this case first,
because it bounds from above what any procedure using the app graph can achieve. We treat
$F \subsetneq \{\, e \in E : e \subseteq A \,\}$ as a later relaxation.

\begin{definition}[True headcount and occupancy]
Each venue holds a definite number of community members at any moment. We call that number its
\emph{true headcount}, written $N_v$. Each venue has a fixed capacity $K_v$, and we call
$N_v / K_v$ its \emph{true occupancy}. The app reports neither quantity.
\end{definition}

The app judges each venue from check-ins posted there in the last two hours, showing every venue
with at least one such check-in. Each shown venue gets one of three levels: quiet, busy, or
packed. The two boundaries are the terciles of tonight's occupancy distribution among open
venues, not a fixed historical threshold. Users compare levels, not numbers.

\begin{definition}[Match rate and coverage]
The app gets a shown venue right when its level matches its true occupancy. The fraction of
shown venues it gets right is the \emph{match rate}. A venue that never appears is a different
kind of error. The fraction of busy-or-packed venues that appear at all is the \emph{coverage}.
\end{definition}

Coverage carries no reference to a procedure. A venue appears when at least one adopter present
posts, whatever the app then computes from the posts. Coverage therefore depends on $C$, on $A$,
and on the posting probabilities, and not on $G$.

\begin{definition}[The question]
Fix a coverage target $\tau \in (0,1)$. We ask which placements of $A$ within $C$ achieve
coverage at least $\tau$.
\end{definition}

We vary structure rather than adoption rate. Two quantities govern the answer: the size of $A$,
and the degree to which $A$ concentrates on $C$. A concentrated $A$ occupies few friend groups.
Those groups fill few venues, and the remaining venues emit nothing. A dispersed $A$ of the same
size meets many groups and misses little. The target $\tau$ therefore determines a curve in those
two quantities, not a single adoption rate.

The app graph serves the levels rather than the coverage. It distinguishes ten check-ins from one
group of friends from ten check-ins by ten strangers, which lets a procedure discount correlated
posts instead of reading them as independent evidence. It also carries information about the
coverage itself, since a dense $G$ indicates that $A$ concentrates on $C$. Its value for the
levels grows with that concentration, because concentration is what generates the correlation it
corrects.

We answer the question inside a synthetic simulation. The simulation generates a venue map, a
flow of simulated people, and a social graph, and it fixes $N_v$ by construction. We use no real
check-in data and no real venue data. We may set the generator's parameters from published
summary statistics, which leaves the graph generated rather than observed.

We compare coverage against the coverage that a uniformly placed adopter set of the same size
achieves. We compare a level assignment against three reference points, fixed before we read any
result: chance, computed rather than assumed; the rule an operator would otherwise ship, meaning
venues ordered by check-in count, by distinct posters, or by count over capacity; and the optimal
rule for this loss, which requires $N_v$ and therefore exists only inside the simulation.

Three limits bound what follows. We report simulation evidence and not a theorem, and the
simulator is our own construction. The concentration of $A$ on $C$ has no external anchor, since
no source we examined reports such a figure for an adopter set on a friendship network, so we
choose the range we sweep. And Fithian, Elith, Hastie and Keith report that thinning leaves
relative intensity unrecoverable when a thinning covariate coincides with an intensity covariate. Assumption 2 lets
posting depend on the venue while occupancy is indexed by the venue, so the two margins of this
problem share a covariate. We treat that result as a soundness check on our level assignments
rather than as a settled objection, because the step transferring it to this setting is unproved.
